\chapter{Threat hunting y por qué grafos}
\label{cap:threat-hunting}

\begin{intuicion}
\emph{Threat hunting} es el oficio de buscar atacantes que ya están
adentro de la red, antes de que detonen. Los \emph{SIEM} clásicos
(\emph{Security Information and Event Management}) son motores de
búsqueda full-text sobre \emph{logs}: andan bien cuando uno sabe qué
buscar (``decime todos los logins de \texttt{admin}''), pero fallan
cuando la señal es \emph{multi-hop}: un usuario inocuo, una máquina
levemente sospechosa, una conexión saliente extraña. Cada evento por
separado no dispara nada; sólo la \emph{cadena} entre ellos delata al
atacante. Modelar el ambiente como un grafo dirigido temporal hace
que esa cadena sea \emph{expresable} en una sola consulta y
\emph{ejecutable} sin recorrer joins manuales.
\end{intuicion}

\section{¿Qué busca un \emph{threat hunter}?}

Imaginá una empresa con 5.000 empleados, 3.000 hosts, una semana de
\emph{logs} de Windows Sysmon, Sentinel y Fortinet. En esa semana
hubo:

\begin{itemize}
    \item $\sim$10 millones de logins exitosos.
    \item $\sim$20 millones de procesos lanzados.
    \item $\sim$80 millones de conexiones de red salientes.
\end{itemize}

\noindent
La pregunta del \emph{threat hunter} no es ``¿quién se logueó hoy?''
(la respuesta sería decenas de miles de filas y ninguna sería
sospechosa por sí misma). La pregunta es:

\begin{quote}
\itshape
``¿Hay algún usuario que se autenticó en un host que después,
en menos de 30~minutos, se conectó a una IP nunca vista antes en
nuestra red, y desde esa IP hubo \emph{beaconing} regular en las
horas siguientes?''
\end{quote}

Esta es una hipótesis \textbf{de tres saltos} (\emph{usuario},
\emph{host}, \emph{IP destino}) con \textbf{restricciones temporales}
(``en menos de 30~minutos'') y \textbf{predicados de comportamiento}
(``IP nunca vista antes'', ``\emph{beaconing} regular''). Cada uno de
los eventos por separado pasa los filtros del SIEM como ruido normal.
La cadena no.

\begin{ejemplo}
Un caso real (en GraphHunter, lo llamamos el ``Sysmon W0'').

\begin{itemize}
    \item \(t_0\): el usuario \texttt{m.gomez} se loguea en
    \texttt{db-01} con credenciales válidas. Login normal: ese
    usuario se loguea acá todos los lunes para sacar reportes.
    \item \(t_0 + 7\,\text{min}\): \texttt{db-01} establece una
    conexión TCP saliente a \texttt{185.159.X.Y:443}. Conexión TLS
    válida, certificado emitido por una CA legítima. La IP no figura
    en la \emph{baseline} de los últimos 30~días.
    \item \(t_0 + 11\,\text{min}\): la conexión se mantiene activa
    con tráfico bajo y regular, $\sim$2~KB cada 60~segundos. Pasa el
    \emph{IDS} sin alertar.
\end{itemize}

Cada uno de esos tres eventos es \emph{individualmente} normal. La
cadena ---usuario válido, host de DB, IP nueva, beacon de bajo
volumen--- es la firma de un \emph{C2} reverse-shell que un atacante
plantó la semana anterior. Para detectarlo, el SIEM tendría que
correr tres queries y \emph{joinar} los resultados a mano.
GraphHunter lo encuentra con una sola hipótesis.
\end{ejemplo}

\section{Por qué un SIEM clásico no alcanza}

Los SIEM operan sobre el modelo \emph{tabla por fuente}: una tabla
para Sysmon, una para Sentinel, una para Fortinet. Cada tabla tiene
sus columnas y sus índices. Una consulta full-text es eficiente sobre
una tabla. Pero ``cadena de tres saltos'' significa
\emph{tres joins} entre tres tablas distintas, donde el atributo que
los une es a veces explícito (\texttt{username}, \texttt{host\_id}) y
a veces implícito (la IP que aparece como \texttt{remote\_ip} en una
conexión saliente debe matchear con la IP que aparece como
\texttt{src\_ip} en el log del firewall).

Tres problemas concretos:

\begin{enumerate}
    \item \textbf{Joins multi-tabla son caros.} Cada \emph{join}
    requiere un índice; el optimizador del SIEM hace lo mejor que
    puede pero la cardinalidad explota. Para tres tablas con 20M
    filas cada una, un \emph{nested-loop join} naive es \(8 \times
    10^{21}\) comparaciones.
    \item \textbf{Filtros temporales son post-join.} El SIEM
    materializa primero el producto cartesiano, después aplica el
    \texttt{WHERE t2 - t1 \(\leq\) 30~min}. Para cuando el filtro
    actúa, el optimizador ya gastó memoria en filas que iban a
    descartarse.
    \item \textbf{La hipótesis no se expresa en una línea.} Escribir
    el join de tres tablas en SQL requiere $\sim$15 líneas, con
    \texttt{JOIN \dots ON \dots AND}, alias, \emph{sub-queries} para
    el ``IP nunca vista antes''. La hipótesis \emph{conceptual} es de
    tres saltos; la representación SQL no lo refleja.
\end{enumerate}

\begin{nota}
Esto no es una crítica al SIEM como producto. Los SIEM están
optimizados para \emph{detección por reglas} (``alertame si veo
\texttt{evt\_id=4625} con más de 5 fallas en 1~minuto'') ---ahí
brillan. Lo que no resuelven bien es la \emph{caza}: encontrar
cadenas que el atacante diseñó precisamente para no disparar reglas
individuales.
\end{nota}

\section{El cambio de modelo: de tablas a grafo temporal}

GraphHunter modela la misma información ---usuarios, hosts, procesos,
IPs, conexiones, autenticaciones--- como un \textbf{grafo dirigido
temporal}:

\begin{itemize}
    \item Los \textbf{nodos} son entidades: \texttt{User},
    \texttt{Host}, \texttt{Process}, \texttt{IP},
    \texttt{File}, \texttt{Domain}.
    \item Las \textbf{aristas} son eventos temporales: ``\texttt{User}
    se autenticó en \texttt{Host} en \(t_1\)'',
    ``\texttt{Host} se conectó a \texttt{IP} en \(t_2\)''.
\end{itemize}

\begin{figure}[H]
\centering
\begin{tikzpicture}[
    node distance=1.4cm,
    every node/.style={font=\small},
    entity/.style={draw=accent, thick, circle, minimum size=10mm,
                   inner sep=1pt, fill=accentlight},
    edge/.style={->, >=Latex, accent, thick}
]
\node[entity] (u) {\strut $u$};
\node[entity, right=2.6cm of u] (h) {\strut $h$};
\node[entity, right=2.6cm of h] (ip) {\strut $\mathit{ip}$};
\node[entity, right=2.6cm of ip] (d) {\strut $d$};

\draw[edge] (u) -- node[above]{\footnotesize\texttt{Auth}\,@\,$t_1$} (h);
\draw[edge] (h) -- node[above]{\footnotesize\texttt{Connect}\,@\,$t_2$} (ip);
\draw[edge] (ip) -- node[above]{\footnotesize\texttt{Beacon}\,@\,$t_3$} (d);

\node[below=1.2cm of h, font=\footnotesize\itshape, accent, align=center]
{hipótesis: \texttt{User -[Auth]-> Host -[Connect]-> IP -[Beacon]-> Domain}\\
 \texttt{WITHIN 30m \{first\_seen\_ip < t1 - 30d\}}};
\end{tikzpicture}
\caption{El ejemplo Sysmon W0 representado como grafo. Cada arista
lleva un \emph{timestamp}; la hipótesis exige
\(t_1 < t_2 < t_3\) y \(t_3 - t_1 \leq 30\)~minutos.}
\label{fig:tx-w0}
\end{figure}

La pregunta de antes ---``¿hay alguna cadena de tres saltos con
estas restricciones?''--- se vuelve, en este modelo,
\emph{enumerar todos los caminos del grafo} que matchean cierta forma
y cumplen ciertos predicados. El motor que lo hace se llama un
\emph{subgraph matching engine}; el capítulo~\ref{cap:subgraph-matching}
lo define con precisión.

\section{Por qué la representación importa}

Tres ventajas concretas del modelo de grafo:

\begin{enumerate}
    \item \textbf{La hipótesis se escribe igual que se piensa.} La
    DSL de GraphHunter (capítulo~\ref{cap:dsl-chain}) deja escribir:
    \begin{lstlisting}
User -[Auth]-> Host -[Connect]-> IP -[Beacon]-> Domain WITHIN 30m
    \end{lstlisting}
    Una línea. La estructura del texto refleja la estructura del
    \emph{path} buscado.

    \item \textbf{El \emph{join} es local.} Recorrer ``vecinos
    salientes de \(v\) con etiqueta \texttt{Connect}'' es \(O(\deg(v))\)
    en una \emph{adjacency list}. No hay producto cartesiano. El
    matcher avanza arista por arista y descarta caminos en cuanto
    violan una restricción temporal o un predicado.

    \item \textbf{Las restricciones temporales son \emph{in-line}.}
    Cuando el matcher elige una arista \((h, \mathit{ip}, t_2)\) para
    extender el \emph{path} parcial \((u, h, t_1)\), chequea
    \(t_1 < t_2\) y \(t_2 - t_1 \leq 30\)~min \emph{antes} de
    recursar. Las ramas que violan la ventana mueren temprano.
\end{enumerate}

\section{Pero modelar como grafo no es magia}

Hay tres preguntas que el resto del libro responde:

\begin{enumerate}
    \item \textbf{¿De dónde sale el grafo?} Un \emph{log} de Sysmon
    es texto JSON; convertirlo a vértices y aristas es la
    \emph{ingesta} (Parte~IV).
    \item \textbf{¿Cómo se almacena un grafo de 100M aristas en una
    notebook con 16~GB de RAM?} Spoiler: con un \texttt{CompactRelation}
    de 32~bytes por arista, \emph{spill} a disco vía mmap, e índices
    auxiliares \emph{lazy} (Partes~III y~VI).
    \item \textbf{¿Cómo se ejecuta el subgraph matching en sub-10~ms
    sobre eso?} Con un zoológico de operadores ---DFS recursivo con
    \emph{prunes}, LFTJ-K\(_4\), Yannakakis chain, Yannakakis
    tree-walk--- y un \emph{planner} que elige el adecuado
    (Parte~V).
\end{enumerate}

\section{Qué viene}

El próximo capítulo formaliza los grafos temporales: definiciones
precisas de vértice, arista, etiqueta, \emph{path}, \emph{walk},
ventana temporal, monotonicidad. El capítulo~\ref{cap:subgraph-matching}
formaliza el problema de \emph{subgraph matching}, distingue
isomorfismo de homomorfismo, y muestra por qué es NP-duro en
general. El capítulo~\ref{cap:joins-y-grafos} construye la conexión
con \emph{joins} relacionales (también desde cero) y deriva el AGM
bound, que es la pieza teórica que justifica los algoritmos
\emph{worst-case-optimal}. El capítulo~\ref{cap:sistemas} cierra la
Parte~I con los fundamentos de sistemas (\emph{cache lines},
\emph{branch prediction}, mmap, atomics) que justifican las
decisiones de \emph{layout} y memoria del motor.

Una vez completados esos cinco capítulos, el resto del libro
construye el motor pieza por pieza, con todas las definiciones ya
disponibles.

\begin{ejercicios}
\begin{enumerate}
    \item \starred{} Escribí en una sola frase ---en castellano, sin
    DSL--- una hipótesis de \emph{lateral movement} de cuatro saltos.
    Identificá los tipos de los vértices y las relaciones entre ellos.

    \item Para la hipótesis del ejemplo Sysmon W0, ¿qué predicados
    tendría que evaluar el motor para descartar usuarios legítimos?
    Listá al menos tres y explicá si cada uno es local al vértice o
    requiere mirar la \emph{baseline} histórica.

    \item \starred{} Considerá un \emph{SIEM} clásico con tres
    tablas de 20M filas cada una. Si el optimizador elige un
    \emph{nested-loop join} naive para una consulta de tres saltos,
    ¿cuántas comparaciones hace en el peor caso? ¿Y si elige un
    \emph{hash join} con cada índice cargado en memoria?

    \item Pensá en una hipótesis multi-hop sobre tu propio dominio
    (web, mobile, e-commerce, IoT) que sea fácil de expresar como
    grafo y difícil de expresar como query SQL. Escribí ambos.

    \item Imaginá que tu grafo tiene \(10^{8}\) aristas. ¿Cuánta
    memoria ocuparía si cada arista ocupara 200 bytes (un
    \texttt{Relation} ``humano'' con strings)? ¿Y si pesara 32 bytes?
    ¿Por qué no se puede almacenar todo el log original y armar el
    grafo \emph{on-demand}?
\end{enumerate}
\end{ejercicios}
